%%%%%%%%%%%%%%%%%%%%%%%%%%%%%%%%%%%%%%%%%%%%%%%%%%%%%%%%%%%%%%%%%%%%%%%%%%
%
%                         华南理工大学开题报告模板
%
%   本文件为模板示例，请根据个人需求进行相应修改。
%
%   1. 请使用 VS Code“打开文件夹”打开“开题报告”目录作为工作空间。
%      本目录的 .vscode/settings.json 会自动生效，无需复制或修改配置。
%      Ctrl+S 保存时默认快速编译；首次使用或更新编号后可手动完整编译。
%   2. 只需编辑本文件中的基本信息、勾选项、摘要和正文。
%      封面、表格、承诺、审核意见和定密审批表由 scutproposal.cls 生成。
%   3. 勾选项留空 {} 时不勾选；填写注释列出的一个值即可勾选对应方框。
%      学位类别在封面和定密审批表中自动同步，无需修改类文件。
%   4. 字体位于上一级 public/fonts/，校名图片位于上一级 public/。
%      编译结果保存到本目录 build/；请保留项目原有目录结构。
%
%%%%%%%%%%%%%%%%%%%%%%%%%%%%%%%%%%%%%%%%%%%%%%%%%%%%%%%%%%%%%%%%%%%%%%%%%%

\documentclass{scutproposal}

%%%%%%%%%%%%%%%%%%%%%%%%%%%%%%%%%%%%%%%%%%%%%%
% 封面及摘要表基本信息
%%%%%%%%%%%%%%%%%%%%%%%%%%%%%%%%%%%%%%%%%%%%%%
% 直接在大括号 {} 里面写就行
\proposalsetup{
  title       = {}, % 中文论文题目
  title-en    = {}, % 英文论文题目
  author      = {}, % 研究生姓名
  student-id  = {}, % 学号
  department  = {}, % 院（系）
  major       = {}, % 专业
  supervisor  = {}, % 指导教师
  date        = {二〇二六年四月}, % 报告日期
  keywords    = {}, % 中文关键词，例如：关键词一 关键词二 关键词三
  keywords-en = {}  % 英文关键词，例如：one / two / three
}

% 勾选项：封面、是否涉密、审核结果及定密审批表，示例如下
% confidential = {no}，表示勾选不涉密
% degree = {master}，表示勾选硕士
\proposalsetup{
  degree         = {}, % master：硕士；doctor：博士
  confidential   = {}, % yes：   涉密；no：不涉密
  review-result  = {}, % pass：  通过；defer：暂缓通过
  employed       = {}, % yes：   在职；no：非在职
  project-source = {}, % government：国家项目；enterprise：委托项目；both：两项都选
  secrecy-level  = {}  % secret：秘密；confidential：机密
}

\begin{document}

\maketitle

% 正文默认顶格，是否缩进由用户决定；需要空两个汉字时，在行首加 \qquad。
% 示例：\qquad 这里填写正文。段落之间留空行；\\ 只换行，不另起段落。

%%%%%%%%%%%%%%%%%%%%%%%%%%%%%%%%%%%%%%%%%%%%%%
% 一、摘要
%%%%%%%%%%%%%%%%%%%%%%%%%%%%%%%%%%%%%%%%%%%%%%
% 不超过 800 字
\begin{proposalabstract}
这是示例摘要
\end{proposalabstract}

%%%%%%%%%%%%%%%%%%%%%%%%%%%%%%%%%%%%%%%%%%%%%%
% 二、立题依据
%%%%%%%%%%%%%%%%%%%%%%%%%%%%%%%%%%%%%%%%%%%%%%
%
\begin{proposalbasis}
\noindent{1、研究意义（工程价值）}

研究意义

\noindent{2、国内外研究现状（工程应用现状）}

研究现状

\noindent{3、主要参考文献及出处}

参考文献

\end{proposalbasis}

%%%%%%%%%%%%%%%%%%%%%%%%%%%%%%%%%%%%%%%%%%%%%%
% 三、研究方案
%%%%%%%%%%%%%%%%%%%%%%%%%%%%%%%%%%%%%%%%%%%%%%
\begin{proposalmethods}
\noindent{1、研究目标、研究内容及拟解决的关键问题}

\noindent{2、拟采取的研究方法、技术路线及可行性分析}

\noindent{3、创新之处}

\end{proposalmethods}

%%%%%%%%%%%%%%%%%%%%%%%%%%%%%%%%%%%%%%%%%%%%%%
% 四、研究条件与基础
%%%%%%%%%%%%%%%%%%%%%%%%%%%%%%%%%%%%%%%%%%%%%%
\begin{proposalfoundation}
\noindent{1、相关的研究工作积累}

\noindent{2、已具备的研究条件，尚缺少的研究条件和拟解决的途径}

\noindent{3、本人已取得的主要成果}

\end{proposalfoundation}

%%%%%%%%%%%%%%%%%%%%%%%%%%%%%%%%%%%%%%%%%%%%%%
% 五、论文工作计划及预期成果
%%%%%%%%%%%%%%%%%%%%%%%%%%%%%%%%%%%%%%%%%%%%%%
% proposalplan 后的花括号填写预期成果，环境内每个 \planitem 对应一行计划。
% 格式：\planitem{起止时间}{工作内容}{备注}；不填备注时保留最后的 {}。
% 未添加计划时自动生成原模板的八行空表；填写时删除示例行前的 %。
\begin{proposalplan}{}
  \planitem{2026年4月--5月}{文献调研与方案设计}{}
  \planitem{}{}{}
  \planitem{}{}{}
  \planitem{}{}{}
  \planitem{}{}{}
  \planitem{}{}{}
  \planitem{}{}{}
  \planitem{}{}{}
\end{proposalplan}

\end{document}
